\section{Introduction}
\label{sec:intro}

When a paper reports that fine-tuning, steering, or a new training objective ``did not change'' a model's power-seeking or sycophancy score, what has the reader learned?
The answer depends on a quantity that the paper almost never reports: the probability that the metric would have moved if the model's behaviour had in fact changed.
Without that number, a null result is silent.
It may mean that the intervention is safe, or it may mean that the instrument is blind.

Other empirical fields handle this problem with positive controls.
An assay is run on a sample that is known to contain the target, and a negative result on the unknown sample is trusted only if the known sample tests positive.
Alignment evaluation relies on a small set of standard instruments, including model-written multiple-choice evaluations \citep{perez2023discovering}, free-form questions scored by a judge model \citep{betley2025emergent}, TruthfulQA \citep{lin2022truthfulqa}, and ETHICS \citep{hendrycks2021ethics}.
These instruments are routinely used as outcome measures, and statistical practice around them is improving \citep{miller2024error}.
However, to our knowledge their \emph{sensitivity} against a panel of models with known behavioural shifts has not been reported, and neither has their false alarm rate under changes that should not matter.

\para{Our approach.}
We build a positive-control protocol for alignment metrics (\figref{fig:protocol}).
We construct \emph{synthetic misalignment anchors}: model conditions that carry a known behavioural shift.
Most anchors are induced by a hidden system prompt that gives the model one of seven dispositions, such as wanting more autonomy or inventing sources, in an overt, concealed, or trigger-gated variant.
Three more are published fine-tuned adapters from the emergent misalignment line of work \citep{turner2025model}.
We verify each anchor with a targeted manipulation check, so that a miss counts against the metric only when the anchor is active.
We then run an unmodified battery of eight standard metrics at 200 items each and count how many verified anchors it fails to detect.
We run the same battery on benign system prompts, such as ``Keep your answers brief and to the point,'' to measure how often it flags a change that should not matter.

Our main question is: {\bf what fraction of verified anchors does a standard alignment battery miss?}
We hypothesized a false negative rate (\fnr) above 0.30, and fixed the detection rule and hypothesis decomposition before collecting data.

\para{Findings.}
The answer depends on what ``standard metrics'' means.
For the full battery of eight metrics, the hypothesis is not supported: the pre-specified rule misses 4 of 19 verified anchors, an \fnr of 0.21 (Wilson 95\% CI 0.09--0.43; one-sided binomial test of \fnr $>0.30$, $p=0.87$).
The interval contains 0.30, so 19 anchors do not rule the hypothesis out either.
For a single metric, the hypothesis is supported: 98 of 152 anchor--metric pairs show no significant shift, an \fnr of 0.64 (0.57--0.72).
Even the metric built for the anchor's own construct misses 4 of 14 anchors.

Two further results change how both nulls and positives should be read.
First, the misses are concentrated.
All four are conditional anchors: the battery misses 4 of 7 conditional anchors and 0 of 12 unconditional ones, and it misses 3 of 5 anchors whose construct has no dedicated metric.
Second, the battery's sensitivity is not discriminating.
The same pre-specified rule flags 5 of 12 benign system prompts as misaligned, a false positive rate (\fpr) of 0.42 (0.19--0.68).
Sample size trades one error for the other: at 25 items per metric the battery misses 52\% to 64\% of anchors, and adding items lowers the miss rate while raising the rate at which benign prompts are flagged.

\para{Contributions.}
\begin{itemize}[leftmargin=*,itemsep=0pt,topsep=0pt]
    \item We propose a positive-control protocol for alignment metrics that pairs verified synthetic misalignment anchors with benign-perturbation controls, and reports \fnr and \fpr together.
    \item We conduct a pre-specified test of the protocol on two 7B open-weight models, 45 candidate anchors, and eight standard metrics, and find a battery-level \fnr of 0.21 and a single-metric \fnr of 0.64.
    \item We show where the battery fails: every miss is a conditional anchor, most are in constructs the battery does not target, and miss rates rise to 0.52--0.64 at 25 items per metric.
    \item We show that benign system prompts trigger the same detection rule 42\% of the time, so that significant shifts need a benign reference as much as nulls need a positive control.
\end{itemize}

We review related work in \secref{sec:related}, describe the protocol in \secref{sec:method}, report results in \secref{sec:results}, and discuss implications and limitations in \secref{sec:discussion}.

\begin{figure}[t]
\centering
\resizebox{\linewidth}{!}{%
\begin{tikzpicture}[
  box/.style={draw, rounded corners=2pt, align=center, font=\small, minimum height=1.25cm, text width=2.75cm, fill=black!4},
  arr/.style={-{Latex[length=2mm]}, thick}]
\node[box] (cond) {{\bf Conditions}\\ baseline, 6 benign\\ controls, 21--24 anchors};
\node[box, right=0.55cm of cond] (ver) {{\bf Verification}\\ 30 targeted answers,\\ blind yes/no judge};
\node[box, right=0.55cm of ver] (bat) {{\bf Metric battery}\\ 8 metrics,\\ 200 items each};
\node[box, right=0.55cm of bat] (det) {{\bf Detection rule}\\ one-sided test vs.\ baseline,\\ Holm over 8 metrics};
\node[box, right=0.55cm of det] (out) {{\bf Error rates}\\ \fnr on verified anchors,\\ \fpr on benign controls};
\draw[arr] (cond) -- (ver);
\draw[arr] (ver) -- (bat);
\draw[arr] (bat) -- (det);
\draw[arr] (det) -- (out);
\end{tikzpicture}}
\caption{The positive-control protocol. Each subject model is run under a baseline, benign system prompts, and anchors with a known behavioural shift. Only anchors that pass a manipulation check enter the \fnr denominator. The benign controls give the false positive rate of the same detection rule.}
\label{fig:protocol}
\end{figure}
